\documentclass[12pt]{article}
\usepackage[margin=1in]{geometry}
\usepackage{enumitem}
\usepackage{titlesec}
\usepackage{parskip}
\usepackage{amsmath}
\usepackage{hyperref}
\usepackage{fontenc}

\title{\textbf{Scott 1985 Claude Guide}\\
\large Scott, James C. \textit{Weapons of the Weak: Everyday Forms of Peasant Resistance}.\\
Bethany, Connecticut: Brevis Press, 1985.}
\author{}
\date{}

\begin{document}

\maketitle
\tableofcontents
\newpage

%--------------------------------------------------
\section{Central Claims}
%--------------------------------------------------

\begin{itemize}[leftmargin=*, itemsep=4pt]
    \item The study of peasant politics has been dominated by attention to organized, collective, confrontational forms of rebellion---uprisings, revolutions, social movements. Scott argues this focus systematically misses the vast bulk of actual peasant political activity, which takes the form of \textbf{everyday resistance}: undramatic, informal, often individual acts that require little or no coordination and avoid direct, symbolic confrontation with authority.
    \item These \textbf{``weapons of the weak''}---foot-dragging, dissimulation, false compliance, feigned ignorance, pilfering, slander, arson, sabotage, desertion---are the ordinary means by which relatively powerless groups defend their interests against elites who extract labor, food, taxes, rent, or interest from them. They require no formal organization, no leadership, no revolutionary consciousness, and can be disavowed if discovered.
    \item Everyday resistance is not merely a poor substitute for rebellion; it is often the more \textbf{rational and effective} strategy, since it minimizes risk to the actor while still imposing real costs on elites (foregone production, damaged property, delayed compliance) and can, in the aggregate, seriously blunt or reshape policies imposed from above---without ever taking the form of an identifiable, punishable collective act.
    \item Class relations in the peasant village are not simply material but also \textbf{discursive/ideological}: resistance is waged as much over the \emph{meaning} of actions and the legitimacy of claims (what counts as ``stinginess,'' ``greed,'' proper charity, or acceptable profit-seeking) as over the material distribution of income. Scott insists on taking this symbolic contest as seriously as the material one.
    \item Scott's ethnographic case---the pseudonymous Malaysian rice-farming village of \textbf{Sedaka} in Kedah state---shows how the \textbf{Green Revolution} (double-cropping via irrigation, and the introduction of combine harvesters) transformed class relations by eliminating the wage labor (transplanting and especially harvesting) on which poor and landless peasants had depended, while enriching an already-advantaged stratum of larger landowners and, at the margins, richer tenant farmers.
    \item Poor peasants respond to this squeeze not with organized revolt (which is largely absent and, given the risks and the fragmented, small-scale nature of the village economy, would be irrational) but with a shifting repertoire of everyday resistance and a vigorous \textbf{ideological campaign}---gossip, character assassination, moral condemnation---aimed at defending the norms of the ``little tradition''/moral economy and stigmatizing those who abandon them.
    \item Scott's account is explicitly pitched against two grand theoretical traditions that, in his view, both miss the texture of actual peasant politics: \textbf{liberal/modernization theory}, which expects passive technological adoption and gradual assimilation into a market economy without serious resistance; and \textbf{orthodox/structural Marxism}, which expects either revolutionary class consciousness or, absent that, ``false consciousness'' produced by ideological hegemony. Scott finds neither: peasants are neither passive nor blind, and their resistance, though rarely revolutionary, is neither mystified nor absent.
    \item The book also foreshadows---without yet fully naming---the distinction Scott develops at length in \textit{Domination and the Arts of Resistance} (1990) between the \textbf{public transcript} (the outwardly deferential, compliant self-presentation peasants offer in the presence of the powerful) and the \textbf{hidden transcript} (the critique, anger, and alternative reading of events articulated offstage, among status equals). Much of Sedaka's everyday resistance and gossip is best understood as the hidden transcript leaking into, or operating just beneath, the public one.
\end{itemize}

%--------------------------------------------------
\section{Core Case Study and Empirical Strategy}
%--------------------------------------------------

\begin{itemize}[leftmargin=*, itemsep=4pt]
    \item \textbf{Sedaka}: a pseudonymous village of roughly seventy households in Kedah, in the rice-growing region of the Muda irrigation scheme in northwest peninsular Malaysia. Scott lived in the village for roughly two years (1978--1980) and combined participant observation with systematic household surveys, life histories, and close attention to gossip, rumor, and informal talk.
    \item \textbf{The Green Revolution as natural experiment}: The Muda irrigation scheme enabled double-cropping (two rice harvests a year instead of one), dramatically raising output. Shortly afterward, combine harvesters were introduced, eliminating the labor-intensive transplanting and (especially) harvesting work that had previously employed poor and landless villagers for wages or in-kind payment (customary harvesting shares).
    \item \textbf{Winners and losers}: Larger landowners captured most of the gains from double-cropping and the labor savings from mechanization; some middle peasants gained by leasing combines or expanding cultivation. Poor peasants and the landless lost their main source of cash and paddy income (harvest wages) and saw the loss of customary gleaning and charitable arrangements tied to hand-harvesting.
    \item \textbf{Methodological strategy}: Rather than studying a revolution, a strike, or an uprising (the conventional objects of peasant-resistance scholarship), Scott deliberately studies a village where no collective, organized resistance occurred--- precisely to demonstrate that the \emph{absence} of dramatic collective action does not mean the absence of political struggle. The struggle is real but takes forms conventional political science overlooks.
    \item \textbf{Sources of evidence}: Detailed household-level data on landholding, tenancy, and income; direct observation of and participation in harvest arrangements; systematic recording of gossip and reputational talk (who is called stingy, who is praised as generous); attention to non-events (rituals of charity not performed, customary payments withheld) as evidence of shifting norms.
    \item \textbf{The village as microcosm}: Sedaka is presented not as a representative sample but as a case that reveals mechanisms---the interplay of economic squeeze, ideological contestation, and small-scale everyday resistance---that Scott argues are generalizable to peasant societies undergoing agrarian change elsewhere, though he is careful to flag Sedaka's particular features (smallholding paddy agriculture, a relatively weak and permeable class structure, absence of overt repression) as scope conditions.
\end{itemize}

%--------------------------------------------------
\section{Core Works Cited}
%--------------------------------------------------

\begin{itemize}[leftmargin=*, itemsep=4pt]
    \item \textbf{Barrington Moore Jr.}, \textit{Social Origins of Dictatorship and Democracy} (1966)---Scott's most important direct interlocutor on peasant revolution. Moore's structural account of the conditions (commercialization of agriculture, strength of the landed elite, character of the state) that produce peasant-based revolution is engaged and revised throughout; Scott argues Moore's framework, built to explain revolutionary outcomes, has nothing to say about the vastly more common non-revolutionary cases.
    \item \textbf{Eric Wolf}, \textit{Peasant Wars of the Twentieth Century} (1969)---another major reference point on the conditions for peasant rebellion (Wolf emphasizes the ``middle peasant'' as the most rebellion-prone stratum); Scott's everyday-resistance framework is presented as filling the enormous empirical space Wolf's revolution-centered account leaves unaddressed.
    \item \textbf{E.P. Thompson}, ``The Moral Economy of the English Crowd in the Eighteenth Century'' (1971) and \textit{Whigs and Hunters} (1975)---foundational for the ``moral economy'' concept and for Thompson's method of reading crowd action, poaching, and ``customary'' resistance as meaningful political behavior with its own logic, not simply criminality or irrationality.
    \item \textbf{Scott's own earlier book}, \textit{The Moral Economy of the Peasant: Rebellion and Subsistence in Southeast Asia} (1976)---develops the subsistence ethic and ``safety-first'' logic of peasant economic behavior that underlies \textit{Weapons of the Weak}'s account of why villagers value charity, reciprocity, and shared poverty norms.
    \item \textbf{Samuel Popkin}, \textit{The Rational Peasant} (1979)---the rival ``political economy''/rational-choice account of peasant behavior that Scott had earlier debated; \textit{Weapons of the Weak} implicitly stakes out a middle position, treating peasants as rational and strategic (contra romantic communitarian readings) while still taking moral-economy norms and ideology seriously (contra pure self-interest accounts).
    \item \textbf{Antonio Gramsci}, \textit{Prison Notebooks} (written 1929--1935)---the source of the hegemony/consent framework that Scott engages critically; Scott argues actual peasant ideological life in Sedaka shows far less genuine consent to elite domination than strong readings of Gramscian hegemony would predict.
    \item \textbf{Clifford Geertz}, \textit{Agricultural Involution} (1963) and \textit{The Religion of Java} (1960)---background ethnographic and interpretive work on Southeast Asian peasant society and the ``little tradition,'' informing Scott's attention to local meaning-systems and religious/normative vocabulary (e.g., \textit{haram}).
    \item \textbf{Michael Adas}, \textit{Prophets of Rebellion} (1979) and work on ``avoidance protest''---parallel scholarship on non-revolutionary peasant responses (flight, millenarianism) to colonial and capitalist pressure, part of the broader turn (with Scott) toward taking non-elite, non-organized political action seriously.
\end{itemize}

%--------------------------------------------------
\section{Chapter 1: The Nature of Peasant Resistance and the Case of Sedaka}
%--------------------------------------------------

\begin{itemize}[leftmargin=*, itemsep=4pt]
    \item \textbf{Critique of the resistance literature}: Scott opens by cataloguing how social science has defined ``real'' resistance so narrowly---organized, collective, principled, revolutionary in intent---that it excludes nearly everything ordinary people actually do to resist domination. This definitional choice, he argues, is not neutral: it reflects elite and state preoccupations with order and rebellion, and it renders peasants either invisible (quiescent) or exceptional (revolutionary).
    \item \textbf{Everyday forms of resistance defined}: ``Foot-dragging, dissimulation, false compliance, feigned ignorance, desertion, pilfering, smuggling, poaching, arson, slander, sabotage, and so on.'' The common features: they require little or no coordination or planning; they represent a form of individual self-help; they typically avoid any direct symbolic confrontation with authority.
    \item \textbf{Why these forms predominate}: Given the risks of open defiance (eviction, loss of wages/tenancy, physical or legal reprisal) and the absence of the organizational resources or protective anonymity that sustained rebellion requires, the rational strategy for the weak is to resist in ways that are deniable and low-risk while still cumulatively effective.
    \item \textbf{Introducing Sedaka}: The village is introduced as a setting for observing this kind of resistance in real time---during and after the disruption of the Green Revolution---rather than as a site of open conflict. Scott is explicit that the absence of collective action in Sedaka should not be mistaken for absence of politics.
    \item \textbf{Class in Sedaka}: A stratified but not sharply polarized peasant community: rich and poor peasants of the same ethnicity, religion, and (frequently) kinship, living in close daily contact---a setting where domination is intimate and continuous, and resistance must operate within ongoing relationships rather than across a clear line of an alien, distant elite.
\end{itemize}

%--------------------------------------------------
\section{Chapter 2: The Ideological Struggle: Stigma and Symbolic Contest}
%--------------------------------------------------

\begin{itemize}[leftmargin=*, itemsep=4pt]
    \item \textbf{Resistance as discourse, not only deed}: Scott insists that class struggle in Sedaka is waged as intensely in words---gossip, character assassination, moral labeling---as in material acts. Whether a rich farmer is called \emph{generous} or \emph{stingy}, whether mechanization is described as \emph{progress} or as \emph{haram} (religiously forbidden/illegitimate), is itself a site of political contest.
    \item \textbf{The vocabulary of stigma}: Poor villagers develop and circulate a shared vocabulary condemning behavior that violates norms of the ``little tradition''---labeling rich farmers who refuse traditional charity, who rent land to outsiders instead of local kin, or who lease out combine harvesters as \emph{kedekut} (stingy), \emph{sombong} (arrogant), or lacking in proper Islamic generosity.
    \item \textbf{Elites' counter-discourse}: Richer villagers respond with their own ideological moves---emphasizing their own hard work and thrift, describing mechanization as inevitable modernization, and reframing charity as a matter of individual discretion rather than binding obligation. The struggle over meaning runs in both directions.
    \item \textbf{Why symbolic contest matters materially}: Reputation has real consequences in a small, face-to-face community---rich peasants who are successfully labeled stingy may lose access to cheap labor, cooperation, or local political support, making the ideological struggle a genuine mechanism of resistance rather than mere talk.
    \item \textbf{Hidden transcript foreshadowed}: Much of this gossip and condemnation happens out of earshot of its targets, among status equals--- an early articulation of what Scott will later call the hidden transcript, contrasted with the outwardly polite, deferential talk (the public transcript) villagers use in face-to-face dealings with the rich.
\end{itemize}

%--------------------------------------------------
\section{Chapter 3: The Green Revolution and the Reshaping of Class Relations}
%--------------------------------------------------

\begin{itemize}[leftmargin=*, itemsep=4pt]
    \item \textbf{Double-cropping}: The Muda irrigation scheme allowed two paddy harvests annually instead of one, roughly doubling potential output and, initially, the demand for labor---a period Scott treats as a relatively good one for poor villagers.
    \item \textbf{Mechanization follows}: The subsequent, rapid adoption of combine harvesters (motivated by rising wage costs and the desire of larger farmers to escape dependence on and social obligations toward local labor) eliminated the bulk of hired harvesting work almost overnight, hitting poor and landless households hardest.
    \item \textbf{Displacement of customary arrangements}: Hand-harvesting had been bound up with customary institutions---the \textit{derau} (mutual harvesting exchange) and religiously sanctioned harvest shares (\textit{zakat peribadi}-adjacent norms)---that combines bypass entirely, severing not just wages but a whole set of reciprocal, charitable relationships between rich and poor.
    \item \textbf{Concentration of gains}: The economic surplus from double-cropping and mechanization accrued overwhelmingly to larger landowners and some middle peasants who could invest in or lease machinery, while poor tenants and laborers saw real income stagnate or decline---a classic case of technical change with sharply unequal, class-patterned distributional effects.
    \item \textbf{Absence of overt conflict}: Despite this squeeze, there is no strike, riot, or organized protest against mechanization in Sedaka---the book's central puzzle, which Scott resolves by redirecting attention to the everyday and ideological forms of resistance detailed in later chapters rather than treating quiescence as proof of consent.
\end{itemize}

%--------------------------------------------------
\section{Chapter 4: The Moral Economy, ``Shared Poverty,'' and the Little Tradition}
%--------------------------------------------------

\begin{itemize}[leftmargin=*, itemsep=4pt]
    \item \textbf{The norm of shared poverty}: Sedaka's ``little tradition'' includes strong normative expectations that the well-off will extend charity, employment, and generosity to poorer kin and neighbors---a moral economy in which subsistence security for the poor is, ideally, a first charge on the community's surplus.
    \item \textbf{Mechanization as a moral, not just material, violation}: Because combine harvesting withdraws work that poor villagers see as customarily and morally theirs, its adoption is experienced and described by the poor not simply as an economic loss but as a \textbf{betrayal of obligation}---rich farmers are seen as reneging on a social contract, not merely making an efficient business decision.
    \item \textbf{Contested legitimacy}: The rich do not accept this framing; they see mechanization as a legitimate exercise of property rights and rational cost-cutting. The moral economy is thus not a settled consensus but itself a contested, actively fought-over set of norms---Scott is careful to avoid romanticizing the little tradition as a harmonious, agreed-upon order.
    \item \textbf{Religion and legitimation}: Islamic vocabulary and obligations (charity, \textit{zakat}, notions of \textit{haram}/\textit{halal}) are mobilized by poorer villagers to frame elite behavior as religiously as well as socially illegitimate, giving the ideological struggle additional moral weight and vocabulary.
    \item \textbf{Continuity with \textit{The Moral Economy of the Peasant} (1976)}: This chapter extends Scott's earlier subsistence-ethic argument, showing the moral economy operating not as a static, precapitalist survival but as a living, contested set of claims actively deployed in response to a very contemporary capitalist transformation (mechanized agriculture).
\end{itemize}

%--------------------------------------------------
\section{Chapter 5: Everyday Forms of Resistance in Practice}
%--------------------------------------------------

\begin{itemize}[leftmargin=*, itemsep=4pt]
    \item \textbf{Concrete repertoire}: Scott catalogues Sedaka's actual resistance practices: foot-dragging on work for rich patrons, minor theft/pilfering of paddy, feigned deference paired with private mockery, refusal of small courtesies, character assassination through gossip, occasional acts of petty sabotage or vandalism, and quiet non-cooperation (declining to attend a rich farmer's ceremony, withholding customary labor).
    \item \textbf{Deniability as design feature}: Each of these tactics is chosen partly because it can be plausibly denied or attributed to accident, laziness, or circumstance rather than intentional defiance---this deniability is what allows resistance to persist without triggering reprisal.
    \item \textbf{Cumulative effect}: Individually trivial, these acts can in aggregate meaningfully affect elite behavior---raising the social and reputational costs of mechanization and stinginess enough to matter, even absent any organized campaign, strike, or petition.
    \item \textbf{Resistance is not solidaristic by default}: Scott is careful to show that everyday resistance among the poor in Sedaka is often uncoordinated, sometimes selfish, occasionally directed against other poor villagers (competition for scarce remaining wage work) rather than a unified class movement---undercutting romantic readings of peasant solidarity.
    \item \textbf{Class as a continuum of practice, not a discrete category}: Because everyday resistance is compatible with continued formal deference and compliance, class relations in Sedaka appear far more muted and ambiguous on the surface than the underlying tension would suggest---exactly the gap between public and hidden transcript that motivates Scott's later work.
\end{itemize}

%--------------------------------------------------
\section{Methodological Contributions and Limitations}
%--------------------------------------------------

\begin{itemize}[leftmargin=*, itemsep=4pt]
    \item \textbf{Long-term ethnography as political science method}: The book is a sustained argument that immersive, extended fieldwork---attentive to gossip, informal talk, and the texture of daily life, not just formal institutions or organized movements---is necessary to detect the actual political behavior of subordinate groups. Survey research or event-based data (strikes, riots, votes) would have found essentially nothing to explain in Sedaka.
    \item \textbf{Combining material and ideational analysis}: Scott's method integrates careful household-level economic data (landholding, tenancy, wage income) with close discourse analysis (the vocabulary of stinginess, generosity, religious legitimacy)---refusing to treat the material and the ideological as separate registers.
    \item \textbf{The falsifiability problem}: A recurring critique (raised by reviewers at the time and since) is that ``everyday resistance,'' defined broadly enough to include foot-dragging, gossip, and feigned ignorance, risks becoming unfalsifiable---nearly any behavior short of enthusiastic compliance can be redescribed as resistance, making it hard to specify what would count as evidence against the thesis.
    \item \textbf{Single-village generalizability}: As with any deep ethnographic case study, questions remain about how far Sedaka's dynamics (a relatively fluid, non-repressive, ethnically homogeneous smallholder community) generalize to peasant societies with sharper class/ethnic boundaries, more repressive states, or different agrarian structures (e.g., plantation economies, sharecropping under debt bondage).
    \item \textbf{Underspecified aggregation mechanism}: Scott demonstrates that everyday resistance can matter but is less precise about exactly how much it constrains elite behavior in Sedaka, or under what conditions everyday resistance escalates into (or forestalls) more organized collective action---a gap later scholarship (including Scott's own subsequent work) tries to address.
    \item \textbf{Positionality and access}: As with much ethnography, the fidelity of the account depends heavily on the fieldworker's access to hidden or backstage talk, particularly among the poor; Scott's extended residence and relationships in the village are the main warrant offered for this access, though it cannot be independently verified in the way survey or archival data can.
\end{itemize}

%--------------------------------------------------
\section{Connections to the Broader Literature}
%--------------------------------------------------

\begin{itemize}[leftmargin=*, itemsep=4pt]
    \item \textbf{Moore 1966}: This is the book's central foil. Moore's \textit{Social Origins of Dictatorship and Democracy} explains variation in regime outcomes (bourgeois democracy, fascism, communist revolution) partly through the political behavior of the peasantry, focusing on the conditions under which peasants become revolutionary allies of one coalition or another. Scott directly engages this framework and argues it is built to explain a rare outcome (peasant revolution) and therefore has almost nothing to say about the overwhelming majority of peasant societies and historical moments in which peasants neither rebel nor passively acquiesce. Where Moore asks ``why do peasants sometimes revolt?'', Scott asks ``what are peasants doing politically the rest of the time?''---a reorientation of the entire research agenda rather than a rival answer to Moore's question.
    \item \textbf{Huntington 1968}: Huntington's concern with political order and the destabilizing effects of rapid modernization without institutionalization parallels Scott's Green Revolution setting (rapid technological/economic change outpacing existing social arrangements), but where Huntington expects instability to manifest as organized unrest or violence when institutions lag behind mobilization, Scott shows that the ``gap'' can instead be absorbed almost entirely into everyday, non-institutional forms of resistance that never rise to the level of a political-order problem in Huntington's sense---a quiet corrective to Huntington's assumption that unmanaged social change produces visible disorder.
    \item \textbf{Dahl 1972}: Dahl's polyarchy framework centers on formal contestation and inclusion within institutionalized channels. Scott's peasants in Sedaka have essentially no access to such channels (no meaningful electoral leverage, no organized interest groups representing the rural poor); \textit{Weapons of the Weak} can be read as an account of political voice and contestation occurring entirely outside the boundaries of Dahl's polyarchical institutions---a reminder that Dahl's framework, built for contexts with functioning contestation, describes only a subset of how less-powerful actors actually exercise voice.
    \item \textbf{Putnam 1993}: Putnam's civic community and social capital operate through voluntary associations and norms of generalized reciprocity that facilitate collective action; Scott's Sedaka has plenty of dense, face-to-face social ties but conspicuously little collective action of the sort Putnam's theory would predict social capital to enable---suggesting that dense local networks can as easily be channeled into individualized, deniable resistance and gossip-based sanctioning as into Putnam-style horizontal civic cooperation. The two books thus offer a useful contrast in what dense social ties are actually used for under conditions of stark inequality versus relative equality.
    \item \textbf{Cox 1997}: Cox's account of how electoral rules shape the incentives for coordination among elites and voters presupposes a functioning electoral arena; Scott's peasants operate almost entirely outside such an arena, which throws into relief how much of Cox's (and similar institutionalist) analysis depends on actors already having formal channels through which to coordinate---channels wholly unavailable to Sedaka's poor.
    \item \textbf{Lijphart 1999}: Lijphart's majoritarian/consensus typology concerns institutional design for aggregating preferences among enfranchised citizens in stable democracies; Scott's subjects are simply outside this frame entirely, which is itself an implicit critique of comparative-institutions scholarship that takes formal democratic participation as the relevant unit of political analysis, ignoring the much larger set of people (agrarian societies, authoritarian contexts, disenfranchised groups) for whom such participation is unavailable.
    \item \textbf{Weber 1930}: Weber's \textit{Protestant Ethic} links religious ethical systems to economic behavior (the Protestant ethic legitimating disciplined accumulation); Scott's Sedaka shows Islamic vocabulary and norms (\textit{haram}, charity obligations) being mobilized in the opposite direction---to delegitimate profit-maximizing behavior (mechanization, withdrawal of charity) rather than to sanctify it---a useful comparative case of religiously inflected economic ethics cutting against, rather than for, capitalist rationalization.
    \item \textbf{Huber and Shipan 2002}: Huber and Shipan's account of policy specificity and bureaucratic discretion concerns principal-agent relationships within formal state institutions; Scott's account of ``foot-dragging'' and ``feigned ignorance'' by the weak toward the powerful is a structurally similar (if informal and asymmetric) principal-agent problem---the poor as reluctant, resistant ``agents'' of elite ``principals''---though operating through social sanction and deniability rather than formal contracts or delegation design.
    \item \textbf{Powell 2000}: Powell's concern is how electoral rules connect citizen preferences to policy outcomes (the chain of representation) in functioning democracies; Scott's peasants have no such chain available to them, so their ``representation'' of interests occurs entirely through non-electoral, informal channels---again illustrating the limits of representation-focused comparative politics for describing the political lives of the rural poor in much of the world.
    \item \textbf{Stokes 2013}: Stokes's work on clientelism and vote-buying examines how patrons secure compliance and loyalty from clients through targeted material exchange within electoral contexts; Scott's Sedaka shows an analogous patron-client material relationship (harvest wages, tenancy, charity) being unilaterally revised by patrons (mechanization) and clients responding not through electoral sanction (unavailable) but through the ideological and everyday-resistance repertoire described above---a non-electoral analogue to the clientelist bargain Stokes studies.
    \item \textbf{Svolik 2012}: Svolik's analysis of authoritarian power-sharing and the constraints rulers face from elites and masses is set at the level of the state; Scott's book operates at the sub-national, village level, but shares an underlying concern with how the relatively powerless constrain the relatively powerful even absent formal institutional checks---informal, reputational, and behavioral constraints substituting for the formal ones Svolik's autocrats also lack vis-à-vis their ruling coalitions.
    \item \textbf{Carey 2009}: Carey's work on legislative voting and party discipline concerns formal institutional mechanisms of compliance and defection within legislatures; Scott's foot-dragging and false compliance are an informal, non-institutional analogue---``defection'' from elite expectations that, like a legislator's vote, can be calibrated to be more or less visible/costly, but here occurs entirely outside any formal institutional record.
    \item \textbf{Kalyvas 2006}: Kalyvas's \textit{The Logic of Violence in Civil War} and Scott's book both center on peasant/local political behavior under conditions of domination and asymmetric power, and both insist on the importance of local, small-scale, often intimate dynamics (rather than macro-level ideology or national politics) in explaining behavior. Where Kalyvas explains why local actors denounce, collaborate, or resist under the shadow of armed actors competing for territorial control (with life-and-death stakes and violence as the key sanctioning mechanism), Scott explains everyday, low-stakes, non-violent resistance under a stable agrarian class structure (with gossip and reputational sanction, not violence, as the key mechanism). Read together, they suggest a spectrum of peasant political behavior under domination, from Scott's low-risk, low-coordination resistance in stable, non-violent settings to Kalyvas's high-risk, identity-revealing choices (collaboration, denunciation, defection) once armed contestation raises the stakes and destroys the ambiguity/deniability that makes Scott's ``weapons of the weak'' effective in the first place.
    \item \textbf{Scott's own later work}: \textit{Domination and the Arts of Resistance} (1990) generalizes and formalizes the public transcript/hidden transcript distinction that is only implicit in \textit{Weapons of the Weak}; \textit{Seeing Like a State} (1998) extends Scott's skepticism of top-down, high-modernist schemes (of which the Green Revolution mechanization drive is an early instance) to a much broader range of state and development projects.
\end{itemize}

\end{document}
